\documentclass[11pt]{article}

\usepackage[a4paper,margin=1in]{geometry}
\usepackage{amsmath,amssymb}
\usepackage{booktabs}
\usepackage{hyperref}
\usepackage{listings}
\usepackage{xcolor}

\lstdefinestyle{cpp}{
  language=C++,
  basicstyle=\ttfamily\small,
  keywordstyle=\color{blue},
  stringstyle=\color{teal},
  commentstyle=\color{gray},
  numbers=left,
  numberstyle=\tiny\color{gray},
  stepnumber=1,
  numbersep=8pt,
  breaklines=true,
  showstringspaces=false,
  tabsize=2
}

\title{Workpiece Convection Boundary Conditions (Thermal Model)}
\author{}
\date{}

\begin{document}
\maketitle
\tableofcontents

\section{Scope}

This document analyzes and documents all heat transfer \emph{convection} properties currently applied to the \emph{workpiece} in the solver's thermal model.

The implementation is not a face-based ``convection boundary condition'' in the classical FEM sense (top/bottom/left/right). Instead, convection is implemented as a \emph{global particle-wise sink term} applied to \emph{all particles} when enabled.

\section{Summary}

\begin{itemize}
  \item Number of convection definitions: 1 (global).
  \item When it is active: cooldown stage only (enabled at the end of cutting if cooldown mode is selected).
  \item Convection coefficient: default $h = 25.0\ \mathrm{W/m^2/K}$ (CLI override).
  \item Ambient fluid temperature: $T_\infty = 295.15\ \mathrm{K}$ (hardcoded as $22^\circ\mathrm{C}$).
  \item Surface assignment: no explicit surface selection; applied to the entire workpiece body (all particles).
  \item Stabilization: global ramp factor to cap maximum predicted cooling rate.
\end{itemize}

\section{Implementation}

\subsection{Where convection is computed}

Convection is computed in the thermal module and applied as part of the thermal RHS assembly:

\begin{itemize}
  \item Convection term: \texttt{src/thermal.cpp} (function \texttt{thermal::apply\_convection}).
  \item Called from: \texttt{thermal::conduction} after conduction discretization.
\end{itemize}

\subsection{Activation and configuration}

Convection is enabled at the start of the cooldown phase in the main program (\texttt{src/refine\_cut\_main.cpp}). Relevant CLI inputs:

\begin{itemize}
  \item \texttt{--cooldown} (enables cooldown stage)
  \item \texttt{--cooldown-hconv <W/m\^{}2/K>} (sets the convection coefficient; also enables cooldown)
\end{itemize}

\section{Parameters}

\subsection{Convection parameters}

\begin{itemize}
  \item Heat transfer coefficient: $h$ in $\mathrm{W/m^2/K}$.
  \item Ambient fluid temperature: $T_\infty$ in $\mathrm{K}$.
  \item Optional rate cap: $\mathrm{max\_rate}$ in $\mathrm{K/s}$.
\end{itemize}

Defaults currently in the program logic:

\begin{itemize}
  \item $h = 25.0\ \mathrm{W/m^2/K}$ (default; CLI override available).
  \item $T_\infty = 273.15 + 22.0 = 295.15\ \mathrm{K}$ (hardcoded).
  \item $\mathrm{max\_rate} = 5.0/60.0\ \mathrm{K/s}$ (cooling-rate cap; hardcoded).
\end{itemize}

\subsection{Particle geometry proxy used for convection}

The code does not compute actual exposed surface area for the workpiece. It computes an \emph{effective} area-to-volume ratio per particle using:

\begin{equation}
dV_i = \frac{m_i}{\rho_0(T_i)}
\end{equation}

\begin{equation}
\left(\frac{A}{V}\right)_i = \frac{4}{\sqrt{dV_i}}
\end{equation}

\section{Applied Convection Properties}

\subsection{BC \#1 --- Cooldown convection (global / particle-wise)}

\subsubsection{Identification}

\begin{itemize}
  \item Name: Cooldown convection (global)
  \item Enabled by: cooldown stage only
  \item Implementation: \texttt{thermal::apply\_convection}
\end{itemize}

\subsubsection{Mathematical form (as coded)}

For each particle $i$:

\begin{equation}
\beta_i = \frac{h\,(A/V)_i}{\rho_0(T_i)\,c_p(T_i)}
\end{equation}

\begin{equation}
\dot{T}_i \mathrel{+}= -r \,\beta_i \,(T_i - T_\infty)
\end{equation}

where:

\begin{itemize}
  \item $h$ is the convection coefficient.
  \item $T_\infty$ is ambient fluid temperature.
  \item $c_p(T)$ is the material heat capacity.
  \item $r$ is a global ramp factor applied when a cooling-rate cap is active.
\end{itemize}

\subsubsection{Cooling-rate ramp (stabilization)}

The solver estimates the maximum predicted convection rate across all particles,

\begin{equation}
\max_i \left|\beta_i (T_i - T_\infty)\right|
\end{equation}

and applies:

\begin{equation}
r=
\begin{cases}
1, & \text{if }\mathrm{max\_rate} \le 0 \text{ or } \max_i|\beta_i (T_i - T_\infty)| \le \mathrm{max\_rate}, \\
\frac{\mathrm{max\_rate}}{\max_i|\beta_i (T_i - T_\infty)|}, & \text{otherwise}.
\end{cases}
\end{equation}

\subsubsection{Surface designation and affected area}

\begin{itemize}
  \item Surface designation: none (no ``top/bottom/side'' mapping).
  \item Affected region: all workpiece particles (including interior).
  \item Area model: implicit per particle via $(A/V)_i = 4/\sqrt{dV_i}$.
\end{itemize}

\subsubsection{Coefficient and ambient temperature}

\begin{center}
\begin{tabular}{lll}
\toprule
Property & Value & Source \\
\midrule
Convection coefficient $h$ & $25.0\ \mathrm{W/m^2\cdot K}$ (default) & \texttt{refine\_cut\_main.cpp} \\
Ambient temperature $T_\infty$ & $295.15\ \mathrm{K}$ ($22^\circ\mathrm{C}$) & \texttt{refine\_cut\_main.cpp} \\
Cooling cap $\mathrm{max\_rate}$ & $0.08333\ \mathrm{K/s}$ ($5^\circ\mathrm{C/min}$) & \texttt{refine\_cut\_main.cpp} \\
\bottomrule
\end{tabular}
\end{center}

\subsubsection{Flow characteristics (velocity/turbulence)}

Not modeled. The code uses a single scalar $h$ without dependence on flow velocity, Reynolds number, turbulence model, or orientation.

\section{Verification Checks}

\subsection{Correct assignment to workpiece surfaces}

The current implementation does not provide surface-based assignment. It applies convection to \emph{all} particles. If the intended operating condition is ``convection only on exposed external boundaries,'' then the current implementation does not match that intent.

\subsection{Missing convection definitions}

\begin{itemize}
  \item No per-surface convection definitions exist (e.g., separate $h$ for top surface vs sides).
  \item No mechanism exists to exclude tool-contact regions or fixed boundaries from convection.
\end{itemize}

\subsection{Duplicate convection definitions}

No duplicates are present. There is a single convection implementation and a single runtime configuration point.

\subsection{Parameter consistency with operating conditions}

\begin{itemize}
  \item Default $h=25\ \mathrm{W/m^2/K}$ is within typical natural convection orders of magnitude for air, but may be inappropriate for forced convection or coolant flows.
  \item $T_\infty$ is hardcoded ($22^\circ\mathrm{C}$), which may not match the intended environment without modifying the source.
\end{itemize}

\section{Findings and Gaps}

\begin{itemize}
  \item The solver's ``convection'' behaves as a global particle-wise Newton cooling sink.
  \item There is no explicit surface selection; ``surface area'' is approximated through a per-particle $(A/V)_i$ proxy.
  \item Convection is not configurable through JSON config; it is enabled and parameterized during cooldown via CLI flags and hardcoded constants.
\end{itemize}

\subsection{Code snippets (for context)}

\begin{lstlisting}[style=cpp]
// src/thermal.cpp (apply_convection)
const double dV = particles[i].m / rho0;
const double A_over_V = 4.0 / std::sqrt(dV);
const double beta = m_h_W_m2K * A_over_V / (rho0 * cp);
particles[i].T_t += -ramp * beta * (Ti - m_T_ambient_K);
\end{lstlisting}

\begin{lstlisting}[style=cpp]
// src/refine_cut_main.cpp (cooldown start)
trml->set_convection(cooldown_hconv_W_m2K, ambient_T_K);
trml->set_max_cooling_rate(max_rate_K_per_s);
trml->set_convection_enabled(true);
\end{lstlisting}

\end{document}

